% introduction
Technological advances have dramatically increased the acquisition of single-cell lineage tracing data, with recently published studies including more than one million cells \citep{yu2026typewriter,colgan2026}. These datasets combine measurements of individual cells with lineage trees describing their division histories and can be generated through live-cell tracking or molecular recording. In live-cell tracking, cell movements are directly observed using time-lapse microscopy, allowing lineage trees to be reconstructed from observed cell divisions \citep{sulston1983lineage,keller2008zebrafish,mcdole2018intoto,ulman2017tracking,ershov2022trackmate}. In molecular recording, cells are engineered to accumulate heritable genomic marks that can subsequently be used to reconstruct their lineage relationships. These marks can be introduced using CRISPR--Cas9 \citep{mckenna2016gestalt,chan2019molecular,bowling2020carlin,frieda2017memoir,chow2021intmemoir}, prime editing \citep{choi2022typewriter,koblan2025}, or base editing \citep{hwang2019l1,winter2026baseline}. Both approaches can pair lineage information with additional cellular measurements, including gene expression \citep{raj2018scgestalt,spanjaard2018linnaeus,bowling2020carlin,colgan2026}, spatial coordinates \citep{koblan2025}, and cell morphology, enabling the study of cell states and behaviors across development, regeneration, and disease.

The scverse ecosystem has established a common framework for single-cell RNA-seq (scRNA-seq) analysis in Python \citep{virshup2023scverse}, centered on AnnData, a standardized data structure for storing single-cell measurements and associated metadata \citep{virshup2024anndata}. Extensions such as MuData and SpatialData add support for multimodal and spatial datasets \citep{bredikhin2022muon,marconato2024spatialdata}. These interoperable data structures support a growing ecosystem of tools, including Scanpy \citep{wolf2018scanpy} and Squidpy \citep{palla2022squidpy}, allowing users to process, visualize, and analyze single-cell data without repeatedly converting between formats. However, the scverse ecosystem has lacked a standardized data structure for lineage trees, making it difficult to integrate lineage analysis into existing single-cell workflows.

Existing tools such as LinTIMaT, DCLEAR, Startle, TiDeTree, and LAML primarily focus on lineage tree reconstruction \citep{zafar2020lintimat,gong2022dclear,sashittal2023startle,seidel2022tidetree,chu2025laml,chu2026lamlpro}. Our Cassiopeia package \citep{jones2020cassiopeia} additionally provides some plotting and analysis functionality, but integrating gene expression data is cumbersome because the \code{CassiopeiaTree} object does not natively store this modality or integrate with the scverse ecosystem. Similarly, users performing lineage analysis with phylogenetic tools such as ETE \citep{huertacepas2016ete} need to manually transfer cell annotations, such as Leiden clusters, from AnnData into the ETE \code{Tree} object. Our goal is to eliminate this data juggling and enable seamless analysis of lineage trees alongside gene expression and other features.

Here, we introduce Pycea, a scverse-compatible toolkit for lineage analysis and visualization built around TreeData, an extension of AnnData that supports the storage of lineage trees alongside other single-cell measurements. TreeData serves as a drop-in replacement for AnnData, retaining compatibility with Scanpy and other scverse tools while enabling integrated lineage analysis. Pycea provides a suite of intuitive and scalable methods for visualizing lineage trees, identifying clades, comparing lineage distances, inferring ancestral states, and more. We demonstrate these capabilities using lineage tracing datasets from mouse development and a metastatic tumor model \citep{colgan2026,koblan2025}. Additional examples and documentation are available in the Pycea documentation \citep{pycea_docs}.
